


% --- Tabla: Caso de Uso ---
\begin{table}[h!]
\centering
\renewcommand{\arraystretch}{1.5}
\begin{tabular}{|l|lllllll|}
\hline
\multicolumn{1}{|l|}{Caso de Uso} & \multicolumn{5}{l|}{\textless{}\textless Nombre del CU \textgreater{}\textgreater{}} & \multicolumn{2}{l|}{\cellcolor[HTML]{E5E5E5}\textless{}\textless Identificador \textgreater{}\textgreater{}} \\ \hline
\multicolumn{1}{|l|}{Actores} & \multicolumn{7}{l|}{\textless{}\textless Listado de los actores participantes en el CU \textgreater{}\textgreater{} \textless{}\textless Podemos indicar quien es el que inicia el CU usando (I) \textgreater{}\textgreater{}} \\ \hline
\multicolumn{1}{|l|}{Tipo} & \multicolumn{7}{l|}{\textless{}\textless Tipo del caso de uso \textgreater{}\textgreater{} \textless{}\textless Primario, Secundario u Opcional | Esencial o Real \textgreater{}\textgreater{}} \\ \hline
\multicolumn{1}{|l|}{Referencias} & \multicolumn{2}{l|}{\textless{}\textless Indicamos qué requisitos se pueden incluir dentro de este CU \textgreater{}\textgreater{}} & \multicolumn{5}{l|}{\textless{}\textless CU que tienen relación con este \textgreater{}\textgreater{}} \\ \hline
\multicolumn{1}{|l|}{Precondición} & \multicolumn{7}{l|}{\textless{}\textless Condiciones sobre el estado del sistema que tienen que ser ciertas para que se pueda realizar el CU \textgreater{}\textgreater{}} \\ \hline
\multicolumn{1}{|l|}{Poscondición} & \multicolumn{7}{l|}{\textless{}\textless Efectos que de forma inmediata tiene la realización del CU sobre el estado del sistema \textgreater{}\textgreater{}} \\ \hline
\multicolumn{1}{|l|}{Autor} & \multicolumn{1}{l|}{\textless{}\textless Esta línea se podría repetir para mantener una historia de cambios del CU \textgreater{}\textgreater{}} & \multicolumn{2}{l|}{Fecha} & \multicolumn{1}{l|}{} & \multicolumn{2}{l|}{Versión} & \multicolumn{1}{l|}{} \\ \hline
\multicolumn{8}{l}{} \\ \hline
\multicolumn{1}{|l|}{\cellcolor[HTML]{F2F2F2}Propósito} & \multicolumn{7}{l|}{\textless{}\textless Descripción general del CU (Suficiente con una línea) \textgreater{}\textgreater{}} \\ \hline
\multicolumn{8}{l}{} \\ \hline
\multicolumn{8}{|l|}{\cellcolor[HTML]{F2F2F2}Resumen} \\ \hline
\multicolumn{8}{|l|}{\textless{}\textless Descripción de alto nivel del flujo normal (básico) del caso de uso (Suficiente con un pequeño párrafo) \textgreater{}\textgreater{}} \\ \hline
\end{tabular}
\end{table}

\vspace{0.8cm}

% --- Tabla: Curso Normal (Básico) ---
\begin{table}[h!]
\centering
\renewcommand{\arraystretch}{1.5}
\begin{tabular}{|p{0.5cm}|p{6cm}|p{6cm}|}
\hline
\multicolumn{3}{|c|}{\cellcolor[HTML]{F2F2F2}\textbf{Curso Normal (Básico)}} \\ 
\hline
1 & \textit{Actor 1: Acción realizada por el actor} & \\
\hline
2 & \textit{Actor 2: Acción realizada por el actor} & \textit{Acción realizada por el sistema} \\ 
\hline
\multirow{2}{*}{N} & \multicolumn{2}{p{12cm}|}{Cuando se realiza la inclusión de otro caso de uso lo representaremos de la forma: \textit{CU-identificador.CU-nombre}} \\ 
\hline
\multicolumn{2}{|p{6.5cm}|}{\textit{<< Se incluyen la secuencia de acciones realizadas por los actores que intervienen en el CU, usando frases cortas que describan el diálogo entre actores y sistema >>}} & \textit{<< Se incluyen las acciones que realiza el sistema ante las acciones de los actores >>} \\
\hline
\end{tabular}
\end{table}

\vspace{0.5cm}

% --- Tabla: Cursos Alternos ---
\begin{table}[h!]
\centering
\renewcommand{\arraystretch}{1.5}
\begin{tabular}{|p{1cm}|p{12cm}|}
\hline
\multicolumn{2}{|c|}{\cellcolor[HTML]{F2F2F2}\textbf{Cursos Alternos}} \\ 
\hline
1a & \textit{Descripción de la secuencia de acciones alternas a la acción 1 del Curso Normal} \\ 
\hline
1b & \textit{<< Secuencia de los cursos alternos del CU >>} \\ 
\hline
\end{tabular}
\end{table}

\vspace{0.5cm}

% --- Tabla: Otros datos ---
\begin{table}[h!]
\centering
\renewcommand{\arraystretch}{1.5}
\begin{tabular}{|p{3cm}|p{4.5cm}|p{3cm}|p{4.5cm}|}
\hline
\multicolumn{4}{|c|}{\cellcolor[HTML]{F2F2F2}\textbf{Otros datos}} \\ 
\hline
\textbf{Frecuencia esperada} & \textit{<< Número de veces que se realiza el CU por unidad de tiempo >>} & \textbf{Rendimiento} & \textit{<< Rendimiento esperado de la secuencia de acciones del CU >>} \\ 
\hline
\textbf{Importancia} & \textit{<< Importancia de este CU en el sistema (vital, alta, moderada, baja) >>} & \textbf{Urgencia} & \textit{<< Urgencia en la realización de este CU durante el desarrollo (alta, moderada, baja) >>} \\ 
\hline
\textbf{Estado} & \textit{<< Estado actual del CU en el desarrollo >>} & \textbf{Estabilidad} & \textit{<< Estabilidad de los requisitos asociados a este CU (alta, moderada, baja) >>} \\ 
\hline
\end{tabular}
\end{table}

\vspace{0.5cm}

% --- Tabla: Comentarios ---
\begin{table}[h!]
\centering
\renewcommand{\arraystretch}{1.5}
\begin{tabular}{|p{15cm}|}
\hline
\cellcolor[HTML]{F2F2F2}\textbf{Comentarios} \\ 
\hline
\textit{<< Comentarios adicionales sobre este CU >>} \\ 
\hline
\end{tabular}
\end{table}
